% GENERATED FILE. Edit canonical lessons, then run npm run generate:books.
% canonical-source-hash: fnv1a64:edbc4ab35c2b4f0a
% canonical-lessons: ML-C496-nadi, ML-C496-aruvi, ML-C496-urava, ML-C496-anakkettu, ML-C496-varruka

\chapter{Rivers and Springs}
\label{ch:ml-rivers-and-springs}
\hlchaptermodality{496}

\begin{chapteropening}
\textbf{By the end of this chapter:} \emph{I can say a large river, a brook, a water spring, a dam and to dry up.}

Complete the last lesson of chapter 496: I can say a large river, a brook, a water spring, a dam and to dry up.
\end{chapteropening}

\section[\texorpdfstring{\ml{നദി} (nadi)}{നദി (nadi)}]{\ml{നദി} (nadi) — a large river}

\label{lesson:ML-C496-nadi}



\begin{quote}
Before the new one: say the Malayalam for a meadow, then the Malayalam for shade.
\end{quote}

\subsection*{You'll want to know: \ml{നദി}}
\textbf{\ml{നദി}} — \emph{nadi} — \textquotedblleft{}a large river\textquotedblright{}.

\begin{quote}
\ml{ആ} \ml{നദി} \ml{കടലിൽ} \ml{ചേരുന്നു}.
\end{quote}

\begin{quote}
\emph{ā nadi kaṭalil cērunnu.}
\end{quote}

\begin{quote}
That river flows into the sea.
\end{quote}

Plural: \ml{നദികൾ}, \emph{nadikaḷ}.

A formal word, used for the big rivers; in speech any river is a \ml{പുഴ}.

\begin{cousinweb}
\textbf{\ml{നദി}} (\emph{nadi}) is borrowed from Sanskrit \emph{nadī}, \textquotedblleft{}a river\textquotedblright{} (Monier-Williams).
\end{cousinweb}

\begin{grammarlens}[title={What you've built}]
One more word for getting around.
\end{grammarlens}

\subsection*{Your turn}
\begin{itemize}
\raggedright
  \item \emph{Say it:} \emph{nadi}
  \item \emph{Say it:} \emph{nadi}, once more
  \item \emph{Say it:} \emph{ā nadi kaṭalil cērunnu.} — That river flows into the sea.
  \item \emph{Recall:} say the Malayalam for a meadow, then the Malayalam for shade, then say \emph{nadi} again
\end{itemize}

\begin{tcolorbox}[breakable,colback=teal!4,colframe=teal!35!black,title={Before you move on}]
What does \ml{നദി} mean? (\textquotedblleft{}A large river\textquotedblright{}.) Say it once more.
\end{tcolorbox}
\section[\texorpdfstring{\ml{അരുവി} (aruvi)}{അരുവി (aruvi)}]{\ml{അരുവി} (aruvi) — a brook; a stream}

\label{lesson:ML-C496-aruvi}



\begin{quote}
Before the new one: say the Malayalam for shade, then the Malayalam for a large river.
\end{quote}

\subsection*{You'll want to know: \ml{അരുവി}}
\textbf{\ml{അരുവി}} — \emph{aruvi} — \textquotedblleft{}a brook; a stream\textquotedblright{}.

\begin{quote}
\ml{കാട്ടിലൂടെ} \ml{ഒരു} \ml{ചെറിയ} \ml{അരുവി} \ml{ഒഴുകുന്നു}.
\end{quote}

\begin{quote}
\emph{kāṭṭilūṭe oru ceṟiya aruvi oḻukunnu.}
\end{quote}

\begin{quote}
A little stream flows through the forest.
\end{quote}

Smaller than a \ml{പുഴ}: water that runs down from the hills.

\begin{cousinweb}
\textbf{\ml{അരുവി}} (\emph{aruvi}) is native Dravidian, from a root that means to spring out and flow down gently (DEDR 226).
\end{cousinweb}

\begin{grammarlens}[title={What you've built}]
One more word for getting around.
\end{grammarlens}

\subsection*{Your turn}
\begin{itemize}
\raggedright
  \item \emph{Say it:} \emph{aruvi}
  \item \emph{Say it:} \emph{aruvi}, once more
  \item \emph{Say it:} \emph{kāṭṭilūṭe oru ceṟiya aruvi oḻukunnu.} — A little stream flows through the forest.
  \item \emph{Recall:} say the Malayalam for shade, then the Malayalam for a large river, then say \emph{aruvi} again
\end{itemize}

\begin{tcolorbox}[breakable,colback=teal!4,colframe=teal!35!black,title={Before you move on}]
What does \ml{അരുവി} mean? (\textquotedblleft{}A brook; a stream\textquotedblright{}.) Say it once more.
\end{tcolorbox}
\section[\texorpdfstring{\ml{ഉറവ} (uṟava)}{ഉറവ (uṟava)}]{\ml{ഉറവ} (uṟava) — a water spring}

\label{lesson:ML-C496-urava}



\begin{quote}
Before the new one: say the Malayalam for a large river, then the Malayalam for a brook.
\end{quote}

\subsection*{You'll want to know: \ml{ഉറവ}}
\textbf{\ml{ഉറവ}} — \emph{uṟava} — \textquotedblleft{}a water spring\textquotedblright{}.

\begin{quote}
\ml{മലയുടെ} \ml{ചുവട്ടിൽ} \ml{ഒരു} \ml{തണുത്ത} \ml{ഉറവ} \ml{ഉണ്ട്}.
\end{quote}

\begin{quote}
\emph{malayuṭe cuvaṭṭil oru taṇutta uṟava uṇṭŭ.}
\end{quote}

\begin{quote}
There is a cold spring at the foot of the hill.
\end{quote}

Where water comes up out of the ground; a well fills from its \ml{ഉറവ}.

\begin{cousinweb}
\textbf{\ml{ഉറവ}} (\emph{uṟava}) is native Dravidian, from a root that means to ooze or spring up (DEDR 761).
\end{cousinweb}

\begin{grammarlens}[title={What you've built}]
One more word for getting around.
\end{grammarlens}

\subsection*{Your turn}
\begin{itemize}
\raggedright
  \item \emph{Say it:} \emph{uṟava}
  \item \emph{Say it:} \emph{uṟava}, once more
  \item \emph{Say it:} \emph{malayuṭe cuvaṭṭil oru taṇutta uṟava uṇṭŭ.} — There is a cold spring at the foot of the hill.
  \item \emph{Recall:} say the Malayalam for a large river, then the Malayalam for a brook, then say \emph{uṟava} again
\end{itemize}

\begin{tcolorbox}[breakable,colback=teal!4,colframe=teal!35!black,title={Before you move on}]
What does \ml{ഉറവ} mean? (\textquotedblleft{}A water spring\textquotedblright{}.) Say it once more.
\end{tcolorbox}
\section[\texorpdfstring{\ml{അണക്കെട്ട്} (aṇakkeṭṭŭ)}{അണക്കെട്ട് (aṇakkeṭṭŭ)}]{\ml{അണക്കെട്ട്} (aṇakkeṭṭŭ) — a dam}

\label{lesson:ML-C496-anakkettu}



\begin{quote}
Before the new one: say the Malayalam for a brook, then the Malayalam for a water spring.
\end{quote}

\subsection*{You'll want to know: \ml{അണക്കെട്ട്}}
\textbf{\ml{അണക്കെട്ട്}} — \emph{aṇakkeṭṭŭ} — \textquotedblleft{}a dam\textquotedblright{}.

\begin{quote}
\ml{മഴക്കാലത്ത്} \ml{അണക്കെട്ടിൽ} \ml{ധാരാളം} \ml{വെള്ളം} \ml{ഉണ്ടാകും}.
\end{quote}

\begin{quote}
\emph{maḻakkālattŭ aṇakkeṭṭil dhārāḷaṁ veḷḷaṁ uṇṭākuṁ.}
\end{quote}

\begin{quote}
In the rainy season the dam holds a lot of water.
\end{quote}

\ml{അണ}, a dam or bank, plus \ml{കെട്ട്}, something built or tied. The shutters are opened when it is full.

\begin{cousinweb}
\textbf{\ml{അണക്കെട്ട്}} (\emph{aṇakkeṭṭŭ}) is a native Dravidian compound of \ml{അണ}, a dam or embankment (DEDR 122), and \ml{കെട്ട്}, a thing built or tied (DEDR 1147).
\end{cousinweb}

\begin{grammarlens}[title={What you've built}]
One more word for getting around.
\end{grammarlens}

\subsection*{Your turn}
\begin{itemize}
\raggedright
  \item \emph{Say it:} \emph{aṇakkeṭṭŭ}
  \item \emph{Say it:} \emph{aṇakkeṭṭŭ}, once more
  \item \emph{Say it:} \emph{maḻakkālattŭ aṇakkeṭṭil dhārāḷaṁ veḷḷaṁ uṇṭākuṁ.} — In the rainy season the dam holds a lot of water.
  \item \emph{Recall:} say the Malayalam for a brook, then the Malayalam for a water spring, then say \emph{aṇakkeṭṭŭ} again
\end{itemize}

\begin{tcolorbox}[breakable,colback=teal!4,colframe=teal!35!black,title={Before you move on}]
What does \ml{അണക്കെട്ട്} mean? (\textquotedblleft{}A dam\textquotedblright{}.) Say it once more.
\end{tcolorbox}
\section[\texorpdfstring{\ml{വറ്റുക} (vaṟṟuka)}{വറ്റുക (vaṟṟuka)}]{\ml{വറ്റുക} (vaṟṟuka) — to dry up}

\label{lesson:ML-C496-varruka}



\begin{quote}
Before the new one: say the Malayalam for a water spring, then the Malayalam for a dam.
\end{quote}

\subsection*{You'll want to know: \ml{വറ്റുക}}
\textbf{\ml{വറ്റുക}} — \emph{vaṟṟuka} — \textquotedblleft{}to dry up\textquotedblright{}.

\begin{quote}
\ml{വേനലിൽ} \ml{ഗ്രാമത്തിലെ} \ml{കിണറുകൾ} \ml{വറ്റി}.
\end{quote}

\begin{quote}
\emph{vēnalil grāmattile kiṇaṟukaḷ vaṟṟi.}
\end{quote}

\begin{quote}
In the dry season the village wells dried up.
\end{quote}

Past: \ml{വറ്റി}, \emph{vaṟṟi}.

Water dries up: a well, a river, a pond. Clothes and grain \ml{ഉണങ്ങുക} instead.

\begin{cousinweb}
\textbf{\ml{വറ്റുക}} (\emph{vaṟṟuka}) is native Dravidian, \textquotedblleft{}to dry up, to grow lean\textquotedblright{} (DEDR 5320).
\end{cousinweb}

\begin{grammarlens}[title={What you've built}]
One more word for getting around.
\end{grammarlens}

\subsection*{Your turn}
\begin{itemize}
\raggedright
  \item \emph{Say it:} \emph{vaṟṟuka}
  \item \emph{Say it:} \emph{vaṟṟuka}, once more
  \item \emph{Say it:} \emph{vēnalil grāmattile kiṇaṟukaḷ vaṟṟi.} — In the dry season the village wells dried up.
  \item \emph{Recall:} say the Malayalam for a water spring, then the Malayalam for a dam, then say \emph{vaṟṟuka} again
\end{itemize}

\begin{tcolorbox}[breakable,colback=teal!4,colframe=teal!35!black,title={Before you move on}]
What does \ml{വറ്റുക} mean? (\textquotedblleft{}To dry up\textquotedblright{}.) Say it once more.
\end{tcolorbox}
